\setcounter{aufgabe}{19}
\hypertarget{e3}{}\einheitenkopf{Einheit 3 von 5 · Punkte und Werte}
\zweigzeile{Hier lernst du, Funktionswerte zu berechnen, Punkte zu prüfen und Nullstellen und Achsenschnittpunkte zu bestimmen · neu in diesem Jahr · P10 oft · baut auf: lineare Funktion (Einheit 2), lineare Gleichungen lösen, negative Zahlen}

\begin{aufgabe}{\textbf{Ich kann eine Zahl für $x$ einsetzen.} Setze die Zahl für $x$ ein und schreibe den Term auf. Rechne noch nicht aus. Negative Zahlen setzt du in Klammern. Beispiel: $f(x) = 4x - 1$, $x = 2$: \quad $f(2) = 4 \cdot 2 - 1$}
\begin{teile}
\teil $f(x) = 5x + 2$, \quad $x = 3$: \quad $f(3) = $ \leerfeld
\teil $f(x) = 2x - 7$, \quad $x = 4$: \quad $f(4) = $ \leerfeld
\teil $f(x) = 3x - 2$, \quad $x = -4$: \quad $f(-4) = $ \leerfeld
\teil $f(x) = -2x + 5$, \quad $x = -3$: \quad $f(-3) = $ \leerfeld
\teil $f(x) = -x + 8$, \quad $x = -6$: \quad $f(-6) = $ \leerfeld
\end{teile}
\end{aufgabe}

\begin{aufgabe}{\textbf{Ich kann Funktionswerte berechnen.} Berechne den Funktionswert.}
\beispiel{f(x) = 4x - 3:\quad f(2) &= 4 \cdot 2 - 3 = 5}
\begin{teilezwei}
\tz $f(x) = 2x + 3$ \quad \feld{f(4)} & \tz $f(x) = 5x - 1$ \quad \feld{f(2)} \\
\tz $f(x) = 3x + 4$ \quad \feld{f(5)} & \tz $f(x) = x + 7$ \quad \feld{f(6)} \\
\tz $f(x) = 3x + 5$ \quad \feld{f(-4)} & \tz $f(x) = 4x - 1$ \quad \feld{f(2{,}5)} \\
\end{teilezwei}
\end{aufgabe}

\begin{aufgabe}{\textbf{Ich kann zu einem Funktionswert das passende $x$ berechnen.} Für welches $x$ hat $f$ den angegebenen Funktionswert? Setze den Term gleich dem Wert und löse die Gleichung.}
\beispiel{f(x) = 3x + 2,\ f(x) = 14:\quad 3x + 2 &= 14 &&\mid -2 \\ 3x &= 12 &&\mid :3 \\ x &= 4}
\begin{gleichungsraster}[3]
\gl{f(x) = 2x + 4,\ f(x) = 10} & \gl{f(x) = 3x - 5,\ f(x) = 7} \\
\gl{f(x) = 5x + 3,\ f(x) = 28} & \gl{f(x) = 4x - 2,\ f(x) = 22} \\
\gl{f(x) = -2x + 6,\ f(x) = 14} & \gl{f(x) = 4x + 7,\ f(x) = 1} \\
\end{gleichungsraster}
\end{aufgabe}

\begin{aufgabe}{\textbf{Ich kann die Nullstelle einer linearen Funktion berechnen.} Setze $f(x) = 0$ und löse nach $x$ auf.}
\beispiel{f(x) = 3x - 6:\quad 0 &= 3x - 6 &&\mid +6 \\ 6 &= 3x &&\mid :3 \\ 2 &= x}
\begin{gleichungsraster}[3]
\gl{f(x) = 2x - 8} & \gl{f(x) = 4x - 12} \\
\gl{f(x) = 5x - 10} & \gl{f(x) = 3x - 21} \\
\end{gleichungsraster}
\end{aufgabe}

\begin{aufgabe}{\textbf{Ich kann die Nullstelle einer linearen Funktion berechnen – weiter.} Berechne die Nullstelle.}
\setcounter{teil}{4}
\begin{gleichungsraster}[3]
\gl{f(x) = x + 3} & \gl{f(x) = -3x + 9} \\
\gl{f(x) = 4x + 6} & \\
\end{gleichungsraster}
\vspace{4mm}
\begin{teile}
\teil Gib die Schnittpunkte des Graphen von $f(x) = -2x + 7$ mit den Achsen an. \\ mit der $x$-Achse: \leerfeld \quad mit der $y$-Achse: \leerfeld
\end{teile}
\end{aufgabe}

\begin{aufgabe}{\textbf{Ich kann prüfen, ob ein Punkt auf einer Geraden liegt.} Prüfe rechnerisch, ob der Punkt auf dem Graphen liegt.}
\beispiel{y = 2x - 3,\ P(4\,|\,5):\quad 5 &= 2 \cdot 4 - 3 \\ 5 &= 5 && \text{wahr: } P \text{ liegt auf dem Graphen}}
\punktprobenliste{x+4/P(2|6), 2x-1/Q(3|6), 4x-2/R(1|3), -x+3/S(5|{-2}), 3x+2/T({-2}|{-4}), -2x+1/U({-3}|{-5})}
\begin{teile}
\teil Liegt der Punkt $V(-8\,|\,{-8})$ auf dem Graphen der Funktion $f$ mit $f(x) = \frac{3}{4}x - 2$? Weise es rechnerisch nach. (P10 2026)
\end{teile}
\end{aufgabe}

\begin{aufgabe}{\textbf{Ich kann eine Wertetabelle ausfüllen und einer Funktion zuordnen.} Fülle die Wertetabelle aus.}
\begin{teile}
\teil $f(x) = x - 3$ \par \wertetabelle{x}{f(x)}{-2,-1,0,1,2}
\teil $f(x) = -2x + 3$ \par \wertetabelle{x}{f(x)}{-2,-1,0,1,2}
\teil $f(x) = 0{,}5x - 1$ \par \wertetabelle{x}{f(x)}{-4,-2,0,2,4}
\end{teile}
Zu welcher Funktion gehört die Tabelle? Kreuze an.
\begin{teile}
\teil \wertetabelle[-4,-1,2,5]{x}{y}{0,1,2,3} \par
\kreuz{$f(x) = 3x - 4$} \\ \kreuz{$f(x) = -4x + 3$} \\ \kreuz{$f(x) = 3x + 4$} \\ \kreuz{$f(x) = x - 4$}
\end{teile}
\end{aufgabe}

\begin{aufgabe}{\textbf{Ich kann Nullstellen und Schnittpunkte am Graphen ablesen.} Lies am Graphen ab.}

\begin{ksys}[xmin=-2,xmax=6,ymin=-5,ymax=6,ablesen]
\gerade[5]{2}{-4}{f}\gerade[-1.5]{-1}{5}{g}
\end{ksys}

\begin{teile}
\teil Nullstelle von $f$: \quad \feld{x}
\teil Schnittpunkt von $f$ mit der $y$-Achse: \quad \leerfeld
\teil Nullstelle von $g$: \quad \feld{x}
\teil Schnittpunkt von $g$ mit der $y$-Achse: \quad \leerfeld
\teil Schnittpunkt der Geraden $f$ und $g$: \quad \leerfeld
\end{teile}
\end{aufgabe}

\begin{aufgabe}{\textbf{Ich finde den Fehler bei der Nullstelle.} Leon soll die Nullstelle von $f(x) = 2x - 6$ bestimmen. Er schreibt:}
\rechnung{f(0) &= 2 \cdot 0 - 6 = -6 && \text{Nullstelle bei } (0\,|\,{-6})}
\begin{teile}
\teil Finde den Fehler, benenne ihn und korrigiere die Rechnung.
\teil Berechne die Nullstelle von $f(x) = 4x - 10$ und gib den Schnittpunkt mit der $y$-Achse an.
\end{teile}
\end{aufgabe}

\begin{aufgabe}{\textbf{Ich kann Aussagen über eine Gerade beurteilen.} Gegeben ist die Funktion $f$ mit $f(x) = -3x + 6$. Ist die Aussage wahr oder falsch? Kreuze an. (P10 2021)}
\begin{teile}
\teil Der Graph von $f$ fällt. \quad \kreuz{wahr}\kreuz{falsch}
\teil Die Nullstelle von $f$ ist 6. \quad \kreuz{wahr}\kreuz{falsch}
\teil Der Punkt $P(-1\,|\,9)$ liegt auf dem Graphen von $f$. \quad \kreuz{wahr}\kreuz{falsch}
\teil Der Graph schneidet die $y$-Achse im Punkt $(0\,|\,{-3})$. \quad \kreuz{wahr}\kreuz{falsch}
\end{teile}
\end{aufgabe}

\begin{aufgabe}{\textbf{Ich kann begründen, wie viele Nullstellen eine Gerade hat.} Begründe.}
\begin{teile}
\teil Warum hat die Funktion $g$ mit $g(x) = 5$ keine Nullstelle?
\teil Warum hat jede lineare Funktion $f(x) = m \cdot x + n$ mit $m \neq 0$ genau eine Nullstelle?
\end{teile}
\end{aufgabe}

\begin{aufgabe}{\textbf{Ich kann Funktionswert und Nullstelle im Sachzusammenhang deuten.} Der Akku eines E-Rollers ist zu 90\,\% geladen. Mit jedem gefahrenen Kilometer sinkt die Ladung um 2,5 Prozentpunkte. Das beschreibt $y = -2{,}5x + 90$ ($x$: Strecke in km, $y$: Ladung in \%).}
\begin{teile}
\teil Wie viel Prozent Ladung hat der Akku nach 12 km?
\teil Nach wie vielen Kilometern sind noch 40\,\% übrig?
\teil Nach wie vielen Kilometern ist der Akku leer? (P10 2021)
\end{teile}
\end{aufgabe}
